\documentclass[10pt,twocolumn]{article}

% =================================================
% arXiv-hardened preamble (pdfLaTeX + natbib + .bbl)
% Technical report, two-column, first-page status bar,
% footer disclaimer, safe package order
% =================================================


\usepackage{authblk}
\renewcommand\Authfont{\large}
\renewcommand\Affilfont{\normalsize}

% -------------------------
% Encoding & Fonts (pdfLaTeX)
% -------------------------
\usepackage[T1]{fontenc}
\usepackage[utf8]{inputenc}
\usepackage{lmodern}
\usepackage{textcomp}

% -------------------------
% Page layout (technical report density)
% -------------------------
\usepackage[a4paper,margin=1.8cm]{geometry}
\setlength{\columnsep}{0.65cm}
\setlength{\emergencystretch}{3em}

% -------------------------
% Spacing
% -------------------------
\usepackage{setspace}
\setstretch{1}
\setlength{\parindent}{0pt}
\setlength{\parskip}{3pt}

% -------------------------
% Section spacing (optional; ok on arXiv)
% -------------------------
\usepackage{titlesec}
\titlespacing*{\section}{0pt}{10pt}{6pt}
\titlespacing*{\subsection}{0pt}{8pt}{4pt}

% -------------------------
% Math
% -------------------------
\usepackage{amsmath,amssymb}

% -------------------------
% Graphics, captions, floats
% -------------------------
\usepackage{graphicx}
\usepackage{xcolor}

% NOTE: Using [H] everywhere in two-column can be painful.
% Keeping float package is fine, but use sparingly.
\usepackage{float}

\usepackage{caption}
\usepackage{subcaption}
\captionsetup[figure]{labelfont=bf, labelsep=colon}
\captionsetup[table]{labelfont=bf}
% Optional: comment this out if floats get stuck
\floatplacement{figure}{H}

% -------------------------
% Tables
% -------------------------
\usepackage{booktabs}
\usepackage{array}
\usepackage{longtable}
\renewcommand{\arraystretch}{1.2}
% NOTE: longtable is not great in twocolumn. Prefer table* / tabular if possible.

% -------------------------
% Icons
% -------------------------
\usepackage{pifont}
\newcommand{\iconbox}[1]{\hspace*{1.6cm}\makebox[1.4em][c]{#1}}
\newcommand{\OK}{\iconbox{\textcolor{green!60!black}{\ding{51}}}}
\newcommand{\WN}{\iconbox{\textcolor{orange!80!black}{\ding{115}}}}
\newcommand{\NK}{\iconbox{\textcolor{red!80!black}{\ding{55}}}}

% -------------------------
% Editorial box
% -------------------------
\usepackage[most]{tcolorbox}
\newtcolorbox{editorialbox}{
  colback=gray!8,
  colframe=gray!50,
  boxrule=0.4pt,
  arc=2pt,
  left=5pt,
  right=5pt,
  top=3pt,
  bottom=3pt
}

% -------------------------
% Footer disclaimer + page number
% -------------------------
\usepackage{fancyhdr}
\pagestyle{fancy}
\fancyhf{}
\fancyfoot[C]{\footnotesize Working Paper -- Not Peer Reviewed}
\fancyfoot[R]{\thepage}
\renewcommand{\headrulewidth}{0pt}

% -------------------------
% Bibliography (natbib; BibTeX; .bbl upload method)
% -------------------------
\usepackage[numbers,sort&compress]{natbib}

% -------------------------
% Hyperlinks (load late)
% -------------------------
\usepackage[hidelinks]{hyperref}
\usepackage{bookmark}  % ok after hyperref
\usepackage{xurl}      % url line breaks
\urlstyle{same}

% -------------------------
% PREPRINT status bar (first page only)
% Put this AFTER hyperref so overlays don't get disturbed by link anchors.
% -------------------------
% \usepackage{tikz}
% \usepackage{eso-pic}

% \newcommand{\statuslabel}{PREPRINT}

% % Draw only on the first page (robust approach):
% \AddToShipoutPictureFG{%
%   \AtPageUpperLeft{%
%     \ifnum\value{page}=1\relax
%       \begin{tikzpicture}[remember picture,overlay]
%         \fill[gray!15]
%           (current page.north west)
%           rectangle
%           ([xshift=1.05cm]current page.south west);

%         \node[rotate=90,anchor=south,
%               gray!60,
%               font=\sffamily\bfseries\footnotesize]
%           at ([xshift=0.525cm,yshift=-0.8cm]current page.north west)
%           {\statuslabel};
%       \end{tikzpicture}%
%     \fi
%   }%
% }

\title{\textbf{Reliable Epistemic AI Tutors for Secondary Education}\\
\large Technical Report}

\author[1]{Wolfgang Spahn}
\author[1]{Marc Eyer}
\affil[1]{PHBern -- University of Teacher Education Bern, Fab 8, Bern, Switzerland}
\affil[ ]{wolfgang.spahn@phbern.ch, marc.eyer@phbern.ch}

\date{February 12, 2026}

\begin{document}

\twocolumn[
\maketitle
\begin{abstract}
The rapid integration of generative Artificial Intelligence (AI) into
educational contexts presents both significant opportunities and
critical risks. While Large Language Models (LLMs) can provide adaptive
feedback and individualized support, they may also foster cognitive
offloading, pedagogical misalignment, and unreliable tutor behavior.
This project introduces \textbf{AIDu}, a reliability-oriented, epistemic
AI tutoring system designed for upper secondary education. AIDu embeds
AI dialogue within authentic, interactive learning environments and
enforces pedagogical role stability through architectural monitoring
mechanisms. Empirical studies with over 180 students demonstrate
significant learning gains and evidence for reducing performance
disparities in abstract STEM domains. The project contributes a scalable
system architecture, a multi-agent AI design, and a formalized approach
to epistemic AI tutoring aligned with regulatory requirements for
high-risk educational AI systems.
\end{abstract}
\vspace{0.8cm}
]

\section{Introduction}\label{introduction}

Generative AI has rapidly become embedded in everyday learning
practices. By 2025, a majority of adolescents reported using AI tools
regularly for academic purposes \citep{tbd}. Yet the empirical
picture remains mixed. Meta-analyses suggest that AI-supported learning
can substantially improve performance \citep{tbd}. At the same time,
experimental evidence on cognitive load and engagement points to
plausible downsides: learners may offload essential sense-making
processes to the system, accept outputs uncritically, and reduce
opportunities for internal knowledge construction and durable
understanding
\citep{kosmyna2025brainchatgptaccumulationcognitive,giannakos2024promise,stadler2024cognitive}.

To mitigate these risks, we argue for educational AI systems that do not
merely provide help, but also \emph{reduce the student's burden of
safeguarding the interaction itself}---for example, by preventing
tutoring-role drift, limiting exposure to hallucinations, and handling
failures in ways that preserve productive learning. General-purpose
conversational agents are not engineered to operate as reliable
pedagogical partners\citep{bang2023education}. They may hallucinate,
shift roles unpredictably, disclose complete solutions prematurely, or
violate basic instructional principles\citep{gan2024llmedu}
\citep{kazemitabaar2024codeaid}. Unlike expert users, secondary
school students cannot be expected to continuously evaluate epistemic
quality, detect subtle errors, or enforce a pedagogically appropriate
framing of the interaction.

Traditional Intelligent Tutoring Systems (ITS) have demonstrated robust
learning gains\citep{anderson1995cognitive}, but often depend on
predefined solution paths and tightly structured domains. While this
supports efficient mastery of well-specified procedures, it can
constrain learner autonomy and may not align with competence-oriented
goals that emphasize transferable skills, strategy use, and
contextualized judgment.

Against this backdrop, this project investigates the following research
question:

\textbf{How must an AI-based tutoring system be designed to reliably
support competence-oriented learning while maintaining pedagogical role
stability and contextual integrity?}

\section{Conceptual Foundation: Epistemic AI
Tutoring}\label{conceptual-foundation-epistemic-ai-tutoring}

AIDu is grounded in the concept of \emph{epistemic tutoring}. In
contrast to assessing, solving, explaining, or strictly Socratic
behaviors, epistemic tutoring emphasizes exploration, hypothesis
generation, reflective feedback, and structural support without
prescribing a fixed solution path.

The design aligns with theories of authentic instruction and competence
development\citep{bransford2000howpeoplelearn},
\citep{newmann1996authentic}, \citep{wood1976scaffolding}.
Learning tasks are structured to require knowledge construction,
higher-order thinking, and transfer beyond classroom contexts. Rather
than reducing cognitive effort, AI support is intended to scaffold
productive struggle. Learners encounter underspecified problems, model
conflicts, and representation mismatches that require active reasoning.

In this framework, AI serves not as a solution provider but as a
structured companion in inquiry processes.

\section{System Architecture}\label{system-architecture}

AIDu is implemented as a modular platform integrating frontend
interfaces, interactive learning environments, backend infrastructure, a
multi-agent AI subsystem, and a monitoring layer for quality assurance.
The architecture is explicitly designed to stabilize pedagogical
behavior and ensure compliance with high-risk AI requirements under the
EU AI Act.

\subsection{Frontend Interfaces}\label{frontend-interfaces}

The browser-based frontend constitutes the primary interaction layer for
learners and teachers. The learner interface enables engagement with
structured learning scenarios, dialogic interaction with the AI tutor,
and exploration of interactive simulations. A functional prototype
allows students to independently work on tasks while receiving real-time
AI feedback. Ongoing developments include a digital whiteboard for
structured solution visualization, formal input tools for mathematical
and chemical expressions, and enhanced coupling between dialogue and
embedded simulations.

The teacher interface provides real-time monitoring of learner activity
and task allocation. Current functionality includes visualization of
system usage states and targeted assignment of learning objectives. A
graphical scenario configuration environment is under development,
enabling teachers to design collaborative tasks and instructional
experiments. Long-term goals include advanced learning analytics
dashboards and individualized trajectory monitoring.

\subsection{Interactive Learning
Spaces}\label{interactive-learning-spaces}

A central element of AIDu is the integration of simulations and digital
experiential environments (``Erfahrungsräume''). These applets enable
active experimentation in domains such as computer science, physics,
mathematics, and foreign language learning. By embedding AI dialogue
within structured simulation contexts, domain boundaries become explicit
and hallucination risks are reduced. The simulations provide dynamic
state information that the AI tutor can reference, ensuring contextual
grounding. Future developments focus on adaptive feedback integration
and personalized learning pathways.

\subsection{Backend Infrastructure}\label{backend-infrastructure}

The backend functions as the coordinating infrastructure of the system.
It manages user authentication, stores learning modules and dialogue
logs, and mediates real-time communication between frontend and AI
components. A scalable reactive architecture supports bidirectional
communication and modular expansion. Administrative management currently
operates via command-line tools; a graphical interface and REST API are
in development to facilitate browser-based configuration. Integration of
over 120 established MINT simulations is supported through an extensible
data model. Planned extensions include support for cooperative learning
scenarios and structured group work.

\subsection{Multi-Agent AI Subsystem}\label{multi-agent-ai-subsystem}

The AI subsystem represents the core innovation of AIDu. It consists of
specialized agents built on LLM architectures. Implemented agents
include an evaluation agent supporting analytical reasoning, a tutoring
agent for STEM domains, a language agent for conversational practice, an
exploration agent for interdisciplinary inquiry, a student simulation
agent for testing, and an experiential space agent coordinating
simulation-based learning.

Agents interact through structured tool calls, enabling them to update
simulation states, annotate the digital whiteboard, extract mathematical
expressions, perform concept analysis, or invoke helper agents.
Mathematical inputs can be formally validated using a Python-based
computer algebra system. For non-mathematical domains, automatic concept
graph generation models relationships among domain concepts and supports
structured feedback.

Currently, interaction flows are partially hard-coded. A dynamic AI
orchestrator is under development to select and combine agents
adaptively based on context and learner behavior. This orchestration
layer will enhance flexibility while preserving pedagogical control.

\subsection{Monitoring and Quality
Assurance}\label{monitoring-and-quality-assurance}

A defining characteristic of AIDu is real-time dialogue monitoring. Each
AI response is evaluated with respect to intended teaching behavior,
session objectives, and domain constraints. Speech-act-based analysis
allows detection of role violations, premature solution delivery, or
misinformation. This reliability layer is essential, as repeated
dialogue simulations demonstrate that identical prompts may otherwise
lead to divergent pedagogical trajectories.

Evaluation procedures include simulated learner testing, mathematical
solution validation, and structured discourse analysis. Planned
extensions involve live dialogue analytics providing teachers with
summaries of interaction depth, tutor behavior patterns, and learner
engagement metrics.

\section{Empirical Evaluation}\label{empirical-evaluation}

The system has been evaluated in multiple studies.

An exploratory study with eight students investigated usability and
dialogue patterns. Results revealed substantial variability in
interaction quality and informed refinements in monitoring mechanisms.

A larger study involving 117 students included structured usability
polling \citep{BuchnerSpahnEyer2026AIDuSGFB}
\citep{Langenegger2025LernenDurchKI} and qualitative analysis of 38
student--AI dialogues. Embedding dialogue within authentic environments
significantly reduced tutor role drift and stabilized epistemic support.

A controlled pretest--posttest study with 84 secondary school students
examined learning outcomes in a networking module focused on diagnosing
a malfunctioning printer. Results showed significant learning gains (d =
0.90). Notably, initial gender performance differences observed in the
pretest disappeared in the posttest, with particularly strong gains
among female students (d = 1.10). These findings suggest that epistemic
AI tutoring embedded in structured environments can both enhance
learning and contribute to reducing disparities in abstract STEM
domains.

Across studies, interaction patterns varied considerably. Some students
used the tutor reflectively for feedback, while others attempted to
extract full solutions. Without architectural safeguards, such pressures
risk undermining epistemic integrity. The monitoring layer proved
essential for maintaining pedagogical alignment.

\section{Contribution and Outlook}\label{contribution-and-outlook}

The AIDu project contributes a formalized model of epistemic AI
tutoring, a reliability-oriented system architecture, and empirical
evidence supporting competence-oriented AI integration in secondary
education. It demonstrates that embedding AI dialogue in authentic
contexts combined with real-time monitoring can stabilize tutor roles
and support reflective learning processes.

The findings underscore that educational impact depends less on raw
model capability than on architectural design. Reliable AI tutors
require explicit pedagogical framing, contextual grounding, dynamic
orchestration, and systematic quality assurance.

Future work will refine adaptive orchestration mechanisms, extend
longitudinal evaluation, and explore scalable deployment in diverse
educational settings. Further research will also investigate cooperative
AI-supported learning and automated real-time dialogue analytics.

\section{Conclusion}\label{conclusion}

AI integration in education is no longer optional. The decisive question
is whether systems are designed to promote durable understanding and
higher-order competencies. AIDu demonstrates that epistemic AI
tutoring---embedded in authentic learning environments and stabilized
through architectural monitoring---can support reflective, self-directed
learning while maintaining pedagogical integrity.

Design, not model size, determines educational reliability.



\appendix
\section*{Appendix A: EU AI Act Alignment and Risk Mitigation}

\subsection*{System Classification}

AIDu qualifies as a high-risk AI system under the European Union AI Act
when deployed in formal educational contexts.

\subsection*{Risk Controls Implemented}

\begin{itemize}
\item Real-time pedagogical role monitoring
\item Detection of epistemic violations and solution leakage
\item Mathematical output validation
\item Contextual grounding via simulation coupling
\item Full interaction logging for auditability
\end{itemize}

\subsection*{Human Oversight}

Teachers retain supervisory authority over system deployment,
task configuration, and intervention in active sessions.

\subsection*{Data Governance}

Interaction data are stored within controlled infrastructure.
The system does not perform automated high-stakes evaluation.

% \clearpage
% \typeout{=== INPUTTING paper.bbl NOW ===}
% \documentclass[10pt,twocolumn]{article}

% =================================================
% arXiv-hardened preamble (pdfLaTeX + natbib + .bbl)
% Technical report, two-column, first-page status bar,
% footer disclaimer, safe package order
% =================================================


\usepackage{authblk}
\renewcommand\Authfont{\large}
\renewcommand\Affilfont{\normalsize}

% -------------------------
% Encoding & Fonts (pdfLaTeX)
% -------------------------
\usepackage[T1]{fontenc}
\usepackage[utf8]{inputenc}
\usepackage{lmodern}
\usepackage{textcomp}

% -------------------------
% Page layout (technical report density)
% -------------------------
\usepackage[a4paper,margin=1.8cm]{geometry}
\setlength{\columnsep}{0.65cm}
\setlength{\emergencystretch}{3em}

% -------------------------
% Spacing
% -------------------------
\usepackage{setspace}
\setstretch{1}
\setlength{\parindent}{0pt}
\setlength{\parskip}{3pt}

% -------------------------
% Section spacing (optional; ok on arXiv)
% -------------------------
\usepackage{titlesec}
\titlespacing*{\section}{0pt}{10pt}{6pt}
\titlespacing*{\subsection}{0pt}{8pt}{4pt}

% -------------------------
% Math
% -------------------------
\usepackage{amsmath,amssymb}

% -------------------------
% Graphics, captions, floats
% -------------------------
\usepackage{graphicx}
\usepackage{xcolor}

% NOTE: Using [H] everywhere in two-column can be painful.
% Keeping float package is fine, but use sparingly.
\usepackage{float}

\usepackage{caption}
\usepackage{subcaption}
\captionsetup[figure]{labelfont=bf, labelsep=colon}
\captionsetup[table]{labelfont=bf}
% Optional: comment this out if floats get stuck
\floatplacement{figure}{H}

% -------------------------
% Tables
% -------------------------
\usepackage{booktabs}
\usepackage{array}
\usepackage{longtable}
\renewcommand{\arraystretch}{1.2}
% NOTE: longtable is not great in twocolumn. Prefer table* / tabular if possible.

% -------------------------
% Icons
% -------------------------
\usepackage{pifont}
\newcommand{\iconbox}[1]{\hspace*{1.6cm}\makebox[1.4em][c]{#1}}
\newcommand{\OK}{\iconbox{\textcolor{green!60!black}{\ding{51}}}}
\newcommand{\WN}{\iconbox{\textcolor{orange!80!black}{\ding{115}}}}
\newcommand{\NK}{\iconbox{\textcolor{red!80!black}{\ding{55}}}}

% -------------------------
% Editorial box
% -------------------------
\usepackage[most]{tcolorbox}
\newtcolorbox{editorialbox}{
  colback=gray!8,
  colframe=gray!50,
  boxrule=0.4pt,
  arc=2pt,
  left=5pt,
  right=5pt,
  top=3pt,
  bottom=3pt
}

% -------------------------
% Footer disclaimer + page number
% -------------------------
\usepackage{fancyhdr}
\pagestyle{fancy}
\fancyhf{}
\fancyfoot[C]{\footnotesize Working Paper -- Not Peer Reviewed}
\fancyfoot[R]{\thepage}
\renewcommand{\headrulewidth}{0pt}

% -------------------------
% Bibliography (natbib; BibTeX; .bbl upload method)
% -------------------------
\usepackage[numbers,sort&compress]{natbib}

% -------------------------
% Hyperlinks (load late)
% -------------------------
\usepackage[hidelinks]{hyperref}
\usepackage{bookmark}  % ok after hyperref
\usepackage{xurl}      % url line breaks
\urlstyle{same}

% -------------------------
% PREPRINT status bar (first page only)
% Put this AFTER hyperref so overlays don't get disturbed by link anchors.
% -------------------------
% \usepackage{tikz}
% \usepackage{eso-pic}

% \newcommand{\statuslabel}{PREPRINT}

% % Draw only on the first page (robust approach):
% \AddToShipoutPictureFG{%
%   \AtPageUpperLeft{%
%     \ifnum\value{page}=1\relax
%       \begin{tikzpicture}[remember picture,overlay]
%         \fill[gray!15]
%           (current page.north west)
%           rectangle
%           ([xshift=1.05cm]current page.south west);

%         \node[rotate=90,anchor=south,
%               gray!60,
%               font=\sffamily\bfseries\footnotesize]
%           at ([xshift=0.525cm,yshift=-0.8cm]current page.north west)
%           {\statuslabel};
%       \end{tikzpicture}%
%     \fi
%   }%
% }

\title{\textbf{Reliable Epistemic AI Tutors for Secondary Education}\\
\large Technical Report}

\author[1]{Wolfgang Spahn}
\author[1]{Marc Eyer}
\affil[1]{PHBern -- University of Teacher Education Bern, Fab 8, Bern, Switzerland}
\affil[ ]{wolfgang.spahn@phbern.ch, marc.eyer@phbern.ch}

\date{February 12, 2026}

\begin{document}

\twocolumn[
\maketitle
\begin{abstract}
The rapid integration of generative Artificial Intelligence (AI) into
educational contexts presents both significant opportunities and
critical risks. While Large Language Models (LLMs) can provide adaptive
feedback and individualized support, they may also foster cognitive
offloading, pedagogical misalignment, and unreliable tutor behavior.
This project introduces \textbf{AIDu}, a reliability-oriented, epistemic
AI tutoring system designed for upper secondary education. AIDu embeds
AI dialogue within authentic, interactive learning environments and
enforces pedagogical role stability through architectural monitoring
mechanisms. Empirical studies with over 180 students demonstrate
significant learning gains and evidence for reducing performance
disparities in abstract STEM domains. The project contributes a scalable
system architecture, a multi-agent AI design, and a formalized approach
to epistemic AI tutoring aligned with regulatory requirements for
high-risk educational AI systems.
\end{abstract}
\vspace{0.8cm}
]

\section{Introduction}\label{introduction}

Generative AI has rapidly become embedded in everyday learning
practices. By 2025, a majority of adolescents reported using AI tools
regularly for academic purposes \citep{tbd}. Yet the empirical
picture remains mixed. Meta-analyses suggest that AI-supported learning
can substantially improve performance \citep{tbd}. At the same time,
experimental evidence on cognitive load and engagement points to
plausible downsides: learners may offload essential sense-making
processes to the system, accept outputs uncritically, and reduce
opportunities for internal knowledge construction and durable
understanding
\citep{kosmyna2025brainchatgptaccumulationcognitive,giannakos2024promise,stadler2024cognitive}.

To mitigate these risks, we argue for educational AI systems that do not
merely provide help, but also \emph{reduce the student's burden of
safeguarding the interaction itself}---for example, by preventing
tutoring-role drift, limiting exposure to hallucinations, and handling
failures in ways that preserve productive learning. General-purpose
conversational agents are not engineered to operate as reliable
pedagogical partners\citep{bang2023education}. They may hallucinate,
shift roles unpredictably, disclose complete solutions prematurely, or
violate basic instructional principles\citep{gan2024llmedu}
\citep{kazemitabaar2024codeaid}. Unlike expert users, secondary
school students cannot be expected to continuously evaluate epistemic
quality, detect subtle errors, or enforce a pedagogically appropriate
framing of the interaction.

Traditional Intelligent Tutoring Systems (ITS) have demonstrated robust
learning gains\citep{anderson1995cognitive}, but often depend on
predefined solution paths and tightly structured domains. While this
supports efficient mastery of well-specified procedures, it can
constrain learner autonomy and may not align with competence-oriented
goals that emphasize transferable skills, strategy use, and
contextualized judgment.

Against this backdrop, this project investigates the following research
question:

\textbf{How must an AI-based tutoring system be designed to reliably
support competence-oriented learning while maintaining pedagogical role
stability and contextual integrity?}

\section{Conceptual Foundation: Epistemic AI
Tutoring}\label{conceptual-foundation-epistemic-ai-tutoring}

AIDu is grounded in the concept of \emph{epistemic tutoring}. In
contrast to assessing, solving, explaining, or strictly Socratic
behaviors, epistemic tutoring emphasizes exploration, hypothesis
generation, reflective feedback, and structural support without
prescribing a fixed solution path.

The design aligns with theories of authentic instruction and competence
development\citep{bransford2000howpeoplelearn},
\citep{newmann1996authentic}, \citep{wood1976scaffolding}.
Learning tasks are structured to require knowledge construction,
higher-order thinking, and transfer beyond classroom contexts. Rather
than reducing cognitive effort, AI support is intended to scaffold
productive struggle. Learners encounter underspecified problems, model
conflicts, and representation mismatches that require active reasoning.

In this framework, AI serves not as a solution provider but as a
structured companion in inquiry processes.

\section{System Architecture}\label{system-architecture}

AIDu is implemented as a modular platform integrating frontend
interfaces, interactive learning environments, backend infrastructure, a
multi-agent AI subsystem, and a monitoring layer for quality assurance.
The architecture is explicitly designed to stabilize pedagogical
behavior and ensure compliance with high-risk AI requirements under the
EU AI Act.

\subsection{Frontend Interfaces}\label{frontend-interfaces}

The browser-based frontend constitutes the primary interaction layer for
learners and teachers. The learner interface enables engagement with
structured learning scenarios, dialogic interaction with the AI tutor,
and exploration of interactive simulations. A functional prototype
allows students to independently work on tasks while receiving real-time
AI feedback. Ongoing developments include a digital whiteboard for
structured solution visualization, formal input tools for mathematical
and chemical expressions, and enhanced coupling between dialogue and
embedded simulations.

The teacher interface provides real-time monitoring of learner activity
and task allocation. Current functionality includes visualization of
system usage states and targeted assignment of learning objectives. A
graphical scenario configuration environment is under development,
enabling teachers to design collaborative tasks and instructional
experiments. Long-term goals include advanced learning analytics
dashboards and individualized trajectory monitoring.

\subsection{Interactive Learning
Spaces}\label{interactive-learning-spaces}

A central element of AIDu is the integration of simulations and digital
experiential environments (``Erfahrungsräume''). These applets enable
active experimentation in domains such as computer science, physics,
mathematics, and foreign language learning. By embedding AI dialogue
within structured simulation contexts, domain boundaries become explicit
and hallucination risks are reduced. The simulations provide dynamic
state information that the AI tutor can reference, ensuring contextual
grounding. Future developments focus on adaptive feedback integration
and personalized learning pathways.

\subsection{Backend Infrastructure}\label{backend-infrastructure}

The backend functions as the coordinating infrastructure of the system.
It manages user authentication, stores learning modules and dialogue
logs, and mediates real-time communication between frontend and AI
components. A scalable reactive architecture supports bidirectional
communication and modular expansion. Administrative management currently
operates via command-line tools; a graphical interface and REST API are
in development to facilitate browser-based configuration. Integration of
over 120 established MINT simulations is supported through an extensible
data model. Planned extensions include support for cooperative learning
scenarios and structured group work.

\subsection{Multi-Agent AI Subsystem}\label{multi-agent-ai-subsystem}

The AI subsystem represents the core innovation of AIDu. It consists of
specialized agents built on LLM architectures. Implemented agents
include an evaluation agent supporting analytical reasoning, a tutoring
agent for STEM domains, a language agent for conversational practice, an
exploration agent for interdisciplinary inquiry, a student simulation
agent for testing, and an experiential space agent coordinating
simulation-based learning.

Agents interact through structured tool calls, enabling them to update
simulation states, annotate the digital whiteboard, extract mathematical
expressions, perform concept analysis, or invoke helper agents.
Mathematical inputs can be formally validated using a Python-based
computer algebra system. For non-mathematical domains, automatic concept
graph generation models relationships among domain concepts and supports
structured feedback.

Currently, interaction flows are partially hard-coded. A dynamic AI
orchestrator is under development to select and combine agents
adaptively based on context and learner behavior. This orchestration
layer will enhance flexibility while preserving pedagogical control.

\subsection{Monitoring and Quality
Assurance}\label{monitoring-and-quality-assurance}

A defining characteristic of AIDu is real-time dialogue monitoring. Each
AI response is evaluated with respect to intended teaching behavior,
session objectives, and domain constraints. Speech-act-based analysis
allows detection of role violations, premature solution delivery, or
misinformation. This reliability layer is essential, as repeated
dialogue simulations demonstrate that identical prompts may otherwise
lead to divergent pedagogical trajectories.

Evaluation procedures include simulated learner testing, mathematical
solution validation, and structured discourse analysis. Planned
extensions involve live dialogue analytics providing teachers with
summaries of interaction depth, tutor behavior patterns, and learner
engagement metrics.

\section{Empirical Evaluation}\label{empirical-evaluation}

The system has been evaluated in multiple studies.

An exploratory study with eight students investigated usability and
dialogue patterns. Results revealed substantial variability in
interaction quality and informed refinements in monitoring mechanisms.

A larger study involving 117 students included structured usability
polling \citep{BuchnerSpahnEyer2026AIDuSGFB}
\citep{Langenegger2025LernenDurchKI} and qualitative analysis of 38
student--AI dialogues. Embedding dialogue within authentic environments
significantly reduced tutor role drift and stabilized epistemic support.

A controlled pretest--posttest study with 84 secondary school students
examined learning outcomes in a networking module focused on diagnosing
a malfunctioning printer. Results showed significant learning gains (d =
0.90). Notably, initial gender performance differences observed in the
pretest disappeared in the posttest, with particularly strong gains
among female students (d = 1.10). These findings suggest that epistemic
AI tutoring embedded in structured environments can both enhance
learning and contribute to reducing disparities in abstract STEM
domains.

Across studies, interaction patterns varied considerably. Some students
used the tutor reflectively for feedback, while others attempted to
extract full solutions. Without architectural safeguards, such pressures
risk undermining epistemic integrity. The monitoring layer proved
essential for maintaining pedagogical alignment.

\section{Contribution and Outlook}\label{contribution-and-outlook}

The AIDu project contributes a formalized model of epistemic AI
tutoring, a reliability-oriented system architecture, and empirical
evidence supporting competence-oriented AI integration in secondary
education. It demonstrates that embedding AI dialogue in authentic
contexts combined with real-time monitoring can stabilize tutor roles
and support reflective learning processes.

The findings underscore that educational impact depends less on raw
model capability than on architectural design. Reliable AI tutors
require explicit pedagogical framing, contextual grounding, dynamic
orchestration, and systematic quality assurance.

Future work will refine adaptive orchestration mechanisms, extend
longitudinal evaluation, and explore scalable deployment in diverse
educational settings. Further research will also investigate cooperative
AI-supported learning and automated real-time dialogue analytics.

\section{Conclusion}\label{conclusion}

AI integration in education is no longer optional. The decisive question
is whether systems are designed to promote durable understanding and
higher-order competencies. AIDu demonstrates that epistemic AI
tutoring---embedded in authentic learning environments and stabilized
through architectural monitoring---can support reflective, self-directed
learning while maintaining pedagogical integrity.

Design, not model size, determines educational reliability.



\appendix
\section*{Appendix A: EU AI Act Alignment and Risk Mitigation}

\subsection*{System Classification}

AIDu qualifies as a high-risk AI system under the European Union AI Act
when deployed in formal educational contexts.

\subsection*{Risk Controls Implemented}

\begin{itemize}
\item Real-time pedagogical role monitoring
\item Detection of epistemic violations and solution leakage
\item Mathematical output validation
\item Contextual grounding via simulation coupling
\item Full interaction logging for auditability
\end{itemize}

\subsection*{Human Oversight}

Teachers retain supervisory authority over system deployment,
task configuration, and intervention in active sessions.

\subsection*{Data Governance}

Interaction data are stored within controlled infrastructure.
The system does not perform automated high-stakes evaluation.

% \clearpage
% \typeout{=== INPUTTING paper.bbl NOW ===}
% \documentclass[10pt,twocolumn]{article}

% =================================================
% arXiv-hardened preamble (pdfLaTeX + natbib + .bbl)
% Technical report, two-column, first-page status bar,
% footer disclaimer, safe package order
% =================================================


\usepackage{authblk}
\renewcommand\Authfont{\large}
\renewcommand\Affilfont{\normalsize}

% -------------------------
% Encoding & Fonts (pdfLaTeX)
% -------------------------
\usepackage[T1]{fontenc}
\usepackage[utf8]{inputenc}
\usepackage{lmodern}
\usepackage{textcomp}

% -------------------------
% Page layout (technical report density)
% -------------------------
\usepackage[a4paper,margin=1.8cm]{geometry}
\setlength{\columnsep}{0.65cm}
\setlength{\emergencystretch}{3em}

% -------------------------
% Spacing
% -------------------------
\usepackage{setspace}
\setstretch{1}
\setlength{\parindent}{0pt}
\setlength{\parskip}{3pt}

% -------------------------
% Section spacing (optional; ok on arXiv)
% -------------------------
\usepackage{titlesec}
\titlespacing*{\section}{0pt}{10pt}{6pt}
\titlespacing*{\subsection}{0pt}{8pt}{4pt}

% -------------------------
% Math
% -------------------------
\usepackage{amsmath,amssymb}

% -------------------------
% Graphics, captions, floats
% -------------------------
\usepackage{graphicx}
\usepackage{xcolor}

% NOTE: Using [H] everywhere in two-column can be painful.
% Keeping float package is fine, but use sparingly.
\usepackage{float}

\usepackage{caption}
\usepackage{subcaption}
\captionsetup[figure]{labelfont=bf, labelsep=colon}
\captionsetup[table]{labelfont=bf}
% Optional: comment this out if floats get stuck
\floatplacement{figure}{H}

% -------------------------
% Tables
% -------------------------
\usepackage{booktabs}
\usepackage{array}
\usepackage{longtable}
\renewcommand{\arraystretch}{1.2}
% NOTE: longtable is not great in twocolumn. Prefer table* / tabular if possible.

% -------------------------
% Icons
% -------------------------
\usepackage{pifont}
\newcommand{\iconbox}[1]{\hspace*{1.6cm}\makebox[1.4em][c]{#1}}
\newcommand{\OK}{\iconbox{\textcolor{green!60!black}{\ding{51}}}}
\newcommand{\WN}{\iconbox{\textcolor{orange!80!black}{\ding{115}}}}
\newcommand{\NK}{\iconbox{\textcolor{red!80!black}{\ding{55}}}}

% -------------------------
% Editorial box
% -------------------------
\usepackage[most]{tcolorbox}
\newtcolorbox{editorialbox}{
  colback=gray!8,
  colframe=gray!50,
  boxrule=0.4pt,
  arc=2pt,
  left=5pt,
  right=5pt,
  top=3pt,
  bottom=3pt
}

% -------------------------
% Footer disclaimer + page number
% -------------------------
\usepackage{fancyhdr}
\pagestyle{fancy}
\fancyhf{}
\fancyfoot[C]{\footnotesize Working Paper -- Not Peer Reviewed}
\fancyfoot[R]{\thepage}
\renewcommand{\headrulewidth}{0pt}

% -------------------------
% Bibliography (natbib; BibTeX; .bbl upload method)
% -------------------------
\usepackage[numbers,sort&compress]{natbib}

% -------------------------
% Hyperlinks (load late)
% -------------------------
\usepackage[hidelinks]{hyperref}
\usepackage{bookmark}  % ok after hyperref
\usepackage{xurl}      % url line breaks
\urlstyle{same}

% -------------------------
% PREPRINT status bar (first page only)
% Put this AFTER hyperref so overlays don't get disturbed by link anchors.
% -------------------------
% \usepackage{tikz}
% \usepackage{eso-pic}

% \newcommand{\statuslabel}{PREPRINT}

% % Draw only on the first page (robust approach):
% \AddToShipoutPictureFG{%
%   \AtPageUpperLeft{%
%     \ifnum\value{page}=1\relax
%       \begin{tikzpicture}[remember picture,overlay]
%         \fill[gray!15]
%           (current page.north west)
%           rectangle
%           ([xshift=1.05cm]current page.south west);

%         \node[rotate=90,anchor=south,
%               gray!60,
%               font=\sffamily\bfseries\footnotesize]
%           at ([xshift=0.525cm,yshift=-0.8cm]current page.north west)
%           {\statuslabel};
%       \end{tikzpicture}%
%     \fi
%   }%
% }

\title{\textbf{Reliable Epistemic AI Tutors for Secondary Education}\\
\large Technical Report}

\author[1]{Wolfgang Spahn}
\author[1]{Marc Eyer}
\affil[1]{PHBern -- University of Teacher Education Bern, Fab 8, Bern, Switzerland}
\affil[ ]{wolfgang.spahn@phbern.ch, marc.eyer@phbern.ch}

\date{February 12, 2026}

\begin{document}

\twocolumn[
\maketitle
\begin{abstract}
The rapid integration of generative Artificial Intelligence (AI) into
educational contexts presents both significant opportunities and
critical risks. While Large Language Models (LLMs) can provide adaptive
feedback and individualized support, they may also foster cognitive
offloading, pedagogical misalignment, and unreliable tutor behavior.
This project introduces \textbf{AIDu}, a reliability-oriented, epistemic
AI tutoring system designed for upper secondary education. AIDu embeds
AI dialogue within authentic, interactive learning environments and
enforces pedagogical role stability through architectural monitoring
mechanisms. Empirical studies with over 180 students demonstrate
significant learning gains and evidence for reducing performance
disparities in abstract STEM domains. The project contributes a scalable
system architecture, a multi-agent AI design, and a formalized approach
to epistemic AI tutoring aligned with regulatory requirements for
high-risk educational AI systems.
\end{abstract}
\vspace{0.8cm}
]

\section{Introduction}\label{introduction}

Generative AI has rapidly become embedded in everyday learning
practices. By 2025, a majority of adolescents reported using AI tools
regularly for academic purposes \citep{tbd}. Yet the empirical
picture remains mixed. Meta-analyses suggest that AI-supported learning
can substantially improve performance \citep{tbd}. At the same time,
experimental evidence on cognitive load and engagement points to
plausible downsides: learners may offload essential sense-making
processes to the system, accept outputs uncritically, and reduce
opportunities for internal knowledge construction and durable
understanding
\citep{kosmyna2025brainchatgptaccumulationcognitive,giannakos2024promise,stadler2024cognitive}.

To mitigate these risks, we argue for educational AI systems that do not
merely provide help, but also \emph{reduce the student's burden of
safeguarding the interaction itself}---for example, by preventing
tutoring-role drift, limiting exposure to hallucinations, and handling
failures in ways that preserve productive learning. General-purpose
conversational agents are not engineered to operate as reliable
pedagogical partners\citep{bang2023education}. They may hallucinate,
shift roles unpredictably, disclose complete solutions prematurely, or
violate basic instructional principles\citep{gan2024llmedu}
\citep{kazemitabaar2024codeaid}. Unlike expert users, secondary
school students cannot be expected to continuously evaluate epistemic
quality, detect subtle errors, or enforce a pedagogically appropriate
framing of the interaction.

Traditional Intelligent Tutoring Systems (ITS) have demonstrated robust
learning gains\citep{anderson1995cognitive}, but often depend on
predefined solution paths and tightly structured domains. While this
supports efficient mastery of well-specified procedures, it can
constrain learner autonomy and may not align with competence-oriented
goals that emphasize transferable skills, strategy use, and
contextualized judgment.

Against this backdrop, this project investigates the following research
question:

\textbf{How must an AI-based tutoring system be designed to reliably
support competence-oriented learning while maintaining pedagogical role
stability and contextual integrity?}

\section{Conceptual Foundation: Epistemic AI
Tutoring}\label{conceptual-foundation-epistemic-ai-tutoring}

AIDu is grounded in the concept of \emph{epistemic tutoring}. In
contrast to assessing, solving, explaining, or strictly Socratic
behaviors, epistemic tutoring emphasizes exploration, hypothesis
generation, reflective feedback, and structural support without
prescribing a fixed solution path.

The design aligns with theories of authentic instruction and competence
development\citep{bransford2000howpeoplelearn},
\citep{newmann1996authentic}, \citep{wood1976scaffolding}.
Learning tasks are structured to require knowledge construction,
higher-order thinking, and transfer beyond classroom contexts. Rather
than reducing cognitive effort, AI support is intended to scaffold
productive struggle. Learners encounter underspecified problems, model
conflicts, and representation mismatches that require active reasoning.

In this framework, AI serves not as a solution provider but as a
structured companion in inquiry processes.

\section{System Architecture}\label{system-architecture}

AIDu is implemented as a modular platform integrating frontend
interfaces, interactive learning environments, backend infrastructure, a
multi-agent AI subsystem, and a monitoring layer for quality assurance.
The architecture is explicitly designed to stabilize pedagogical
behavior and ensure compliance with high-risk AI requirements under the
EU AI Act.

\subsection{Frontend Interfaces}\label{frontend-interfaces}

The browser-based frontend constitutes the primary interaction layer for
learners and teachers. The learner interface enables engagement with
structured learning scenarios, dialogic interaction with the AI tutor,
and exploration of interactive simulations. A functional prototype
allows students to independently work on tasks while receiving real-time
AI feedback. Ongoing developments include a digital whiteboard for
structured solution visualization, formal input tools for mathematical
and chemical expressions, and enhanced coupling between dialogue and
embedded simulations.

The teacher interface provides real-time monitoring of learner activity
and task allocation. Current functionality includes visualization of
system usage states and targeted assignment of learning objectives. A
graphical scenario configuration environment is under development,
enabling teachers to design collaborative tasks and instructional
experiments. Long-term goals include advanced learning analytics
dashboards and individualized trajectory monitoring.

\subsection{Interactive Learning
Spaces}\label{interactive-learning-spaces}

A central element of AIDu is the integration of simulations and digital
experiential environments (``Erfahrungsräume''). These applets enable
active experimentation in domains such as computer science, physics,
mathematics, and foreign language learning. By embedding AI dialogue
within structured simulation contexts, domain boundaries become explicit
and hallucination risks are reduced. The simulations provide dynamic
state information that the AI tutor can reference, ensuring contextual
grounding. Future developments focus on adaptive feedback integration
and personalized learning pathways.

\subsection{Backend Infrastructure}\label{backend-infrastructure}

The backend functions as the coordinating infrastructure of the system.
It manages user authentication, stores learning modules and dialogue
logs, and mediates real-time communication between frontend and AI
components. A scalable reactive architecture supports bidirectional
communication and modular expansion. Administrative management currently
operates via command-line tools; a graphical interface and REST API are
in development to facilitate browser-based configuration. Integration of
over 120 established MINT simulations is supported through an extensible
data model. Planned extensions include support for cooperative learning
scenarios and structured group work.

\subsection{Multi-Agent AI Subsystem}\label{multi-agent-ai-subsystem}

The AI subsystem represents the core innovation of AIDu. It consists of
specialized agents built on LLM architectures. Implemented agents
include an evaluation agent supporting analytical reasoning, a tutoring
agent for STEM domains, a language agent for conversational practice, an
exploration agent for interdisciplinary inquiry, a student simulation
agent for testing, and an experiential space agent coordinating
simulation-based learning.

Agents interact through structured tool calls, enabling them to update
simulation states, annotate the digital whiteboard, extract mathematical
expressions, perform concept analysis, or invoke helper agents.
Mathematical inputs can be formally validated using a Python-based
computer algebra system. For non-mathematical domains, automatic concept
graph generation models relationships among domain concepts and supports
structured feedback.

Currently, interaction flows are partially hard-coded. A dynamic AI
orchestrator is under development to select and combine agents
adaptively based on context and learner behavior. This orchestration
layer will enhance flexibility while preserving pedagogical control.

\subsection{Monitoring and Quality
Assurance}\label{monitoring-and-quality-assurance}

A defining characteristic of AIDu is real-time dialogue monitoring. Each
AI response is evaluated with respect to intended teaching behavior,
session objectives, and domain constraints. Speech-act-based analysis
allows detection of role violations, premature solution delivery, or
misinformation. This reliability layer is essential, as repeated
dialogue simulations demonstrate that identical prompts may otherwise
lead to divergent pedagogical trajectories.

Evaluation procedures include simulated learner testing, mathematical
solution validation, and structured discourse analysis. Planned
extensions involve live dialogue analytics providing teachers with
summaries of interaction depth, tutor behavior patterns, and learner
engagement metrics.

\section{Empirical Evaluation}\label{empirical-evaluation}

The system has been evaluated in multiple studies.

An exploratory study with eight students investigated usability and
dialogue patterns. Results revealed substantial variability in
interaction quality and informed refinements in monitoring mechanisms.

A larger study involving 117 students included structured usability
polling \citep{BuchnerSpahnEyer2026AIDuSGFB}
\citep{Langenegger2025LernenDurchKI} and qualitative analysis of 38
student--AI dialogues. Embedding dialogue within authentic environments
significantly reduced tutor role drift and stabilized epistemic support.

A controlled pretest--posttest study with 84 secondary school students
examined learning outcomes in a networking module focused on diagnosing
a malfunctioning printer. Results showed significant learning gains (d =
0.90). Notably, initial gender performance differences observed in the
pretest disappeared in the posttest, with particularly strong gains
among female students (d = 1.10). These findings suggest that epistemic
AI tutoring embedded in structured environments can both enhance
learning and contribute to reducing disparities in abstract STEM
domains.

Across studies, interaction patterns varied considerably. Some students
used the tutor reflectively for feedback, while others attempted to
extract full solutions. Without architectural safeguards, such pressures
risk undermining epistemic integrity. The monitoring layer proved
essential for maintaining pedagogical alignment.

\section{Contribution and Outlook}\label{contribution-and-outlook}

The AIDu project contributes a formalized model of epistemic AI
tutoring, a reliability-oriented system architecture, and empirical
evidence supporting competence-oriented AI integration in secondary
education. It demonstrates that embedding AI dialogue in authentic
contexts combined with real-time monitoring can stabilize tutor roles
and support reflective learning processes.

The findings underscore that educational impact depends less on raw
model capability than on architectural design. Reliable AI tutors
require explicit pedagogical framing, contextual grounding, dynamic
orchestration, and systematic quality assurance.

Future work will refine adaptive orchestration mechanisms, extend
longitudinal evaluation, and explore scalable deployment in diverse
educational settings. Further research will also investigate cooperative
AI-supported learning and automated real-time dialogue analytics.

\section{Conclusion}\label{conclusion}

AI integration in education is no longer optional. The decisive question
is whether systems are designed to promote durable understanding and
higher-order competencies. AIDu demonstrates that epistemic AI
tutoring---embedded in authentic learning environments and stabilized
through architectural monitoring---can support reflective, self-directed
learning while maintaining pedagogical integrity.

Design, not model size, determines educational reliability.



\appendix
\section*{Appendix A: EU AI Act Alignment and Risk Mitigation}

\subsection*{System Classification}

AIDu qualifies as a high-risk AI system under the European Union AI Act
when deployed in formal educational contexts.

\subsection*{Risk Controls Implemented}

\begin{itemize}
\item Real-time pedagogical role monitoring
\item Detection of epistemic violations and solution leakage
\item Mathematical output validation
\item Contextual grounding via simulation coupling
\item Full interaction logging for auditability
\end{itemize}

\subsection*{Human Oversight}

Teachers retain supervisory authority over system deployment,
task configuration, and intervention in active sessions.

\subsection*{Data Governance}

Interaction data are stored within controlled infrastructure.
The system does not perform automated high-stakes evaluation.

% \clearpage
% \typeout{=== INPUTTING paper.bbl NOW ===}
% \documentclass[10pt,twocolumn]{article}

\input{preamble.tex}

\title{\textbf{Reliable Epistemic AI Tutors for Secondary Education}\\
\large Technical Report}

\author[1]{Wolfgang Spahn}
\author[1]{Marc Eyer}
\affil[1]{PHBern -- University of Teacher Education Bern, Fab 8, Bern, Switzerland}
\affil[ ]{wolfgang.spahn@phbern.ch, marc.eyer@phbern.ch}

\date{February 12, 2026}

\begin{document}

\twocolumn[
\maketitle
\begin{abstract}
The rapid integration of generative Artificial Intelligence (AI) into
educational contexts presents both significant opportunities and
critical risks. While Large Language Models (LLMs) can provide adaptive
feedback and individualized support, they may also foster cognitive
offloading, pedagogical misalignment, and unreliable tutor behavior.
This project introduces \textbf{AIDu}, a reliability-oriented, epistemic
AI tutoring system designed for upper secondary education. AIDu embeds
AI dialogue within authentic, interactive learning environments and
enforces pedagogical role stability through architectural monitoring
mechanisms. Empirical studies with over 180 students demonstrate
significant learning gains and evidence for reducing performance
disparities in abstract STEM domains. The project contributes a scalable
system architecture, a multi-agent AI design, and a formalized approach
to epistemic AI tutoring aligned with regulatory requirements for
high-risk educational AI systems.
\end{abstract}
\vspace{0.8cm}
]

\section{Introduction}\label{introduction}

Generative AI has rapidly become embedded in everyday learning
practices. By 2025, a majority of adolescents reported using AI tools
regularly for academic purposes \citep{tbd}. Yet the empirical
picture remains mixed. Meta-analyses suggest that AI-supported learning
can substantially improve performance \citep{tbd}. At the same time,
experimental evidence on cognitive load and engagement points to
plausible downsides: learners may offload essential sense-making
processes to the system, accept outputs uncritically, and reduce
opportunities for internal knowledge construction and durable
understanding
\citep{kosmyna2025brainchatgptaccumulationcognitive,giannakos2024promise,stadler2024cognitive}.

To mitigate these risks, we argue for educational AI systems that do not
merely provide help, but also \emph{reduce the student's burden of
safeguarding the interaction itself}---for example, by preventing
tutoring-role drift, limiting exposure to hallucinations, and handling
failures in ways that preserve productive learning. General-purpose
conversational agents are not engineered to operate as reliable
pedagogical partners\citep{bang2023education}. They may hallucinate,
shift roles unpredictably, disclose complete solutions prematurely, or
violate basic instructional principles\citep{gan2024llmedu}
\citep{kazemitabaar2024codeaid}. Unlike expert users, secondary
school students cannot be expected to continuously evaluate epistemic
quality, detect subtle errors, or enforce a pedagogically appropriate
framing of the interaction.

Traditional Intelligent Tutoring Systems (ITS) have demonstrated robust
learning gains\citep{anderson1995cognitive}, but often depend on
predefined solution paths and tightly structured domains. While this
supports efficient mastery of well-specified procedures, it can
constrain learner autonomy and may not align with competence-oriented
goals that emphasize transferable skills, strategy use, and
contextualized judgment.

Against this backdrop, this project investigates the following research
question:

\textbf{How must an AI-based tutoring system be designed to reliably
support competence-oriented learning while maintaining pedagogical role
stability and contextual integrity?}

\section{Conceptual Foundation: Epistemic AI
Tutoring}\label{conceptual-foundation-epistemic-ai-tutoring}

AIDu is grounded in the concept of \emph{epistemic tutoring}. In
contrast to assessing, solving, explaining, or strictly Socratic
behaviors, epistemic tutoring emphasizes exploration, hypothesis
generation, reflective feedback, and structural support without
prescribing a fixed solution path.

The design aligns with theories of authentic instruction and competence
development\citep{bransford2000howpeoplelearn},
\citep{newmann1996authentic}, \citep{wood1976scaffolding}.
Learning tasks are structured to require knowledge construction,
higher-order thinking, and transfer beyond classroom contexts. Rather
than reducing cognitive effort, AI support is intended to scaffold
productive struggle. Learners encounter underspecified problems, model
conflicts, and representation mismatches that require active reasoning.

In this framework, AI serves not as a solution provider but as a
structured companion in inquiry processes.

\section{System Architecture}\label{system-architecture}

AIDu is implemented as a modular platform integrating frontend
interfaces, interactive learning environments, backend infrastructure, a
multi-agent AI subsystem, and a monitoring layer for quality assurance.
The architecture is explicitly designed to stabilize pedagogical
behavior and ensure compliance with high-risk AI requirements under the
EU AI Act.

\subsection{Frontend Interfaces}\label{frontend-interfaces}

The browser-based frontend constitutes the primary interaction layer for
learners and teachers. The learner interface enables engagement with
structured learning scenarios, dialogic interaction with the AI tutor,
and exploration of interactive simulations. A functional prototype
allows students to independently work on tasks while receiving real-time
AI feedback. Ongoing developments include a digital whiteboard for
structured solution visualization, formal input tools for mathematical
and chemical expressions, and enhanced coupling between dialogue and
embedded simulations.

The teacher interface provides real-time monitoring of learner activity
and task allocation. Current functionality includes visualization of
system usage states and targeted assignment of learning objectives. A
graphical scenario configuration environment is under development,
enabling teachers to design collaborative tasks and instructional
experiments. Long-term goals include advanced learning analytics
dashboards and individualized trajectory monitoring.

\subsection{Interactive Learning
Spaces}\label{interactive-learning-spaces}

A central element of AIDu is the integration of simulations and digital
experiential environments (``Erfahrungsräume''). These applets enable
active experimentation in domains such as computer science, physics,
mathematics, and foreign language learning. By embedding AI dialogue
within structured simulation contexts, domain boundaries become explicit
and hallucination risks are reduced. The simulations provide dynamic
state information that the AI tutor can reference, ensuring contextual
grounding. Future developments focus on adaptive feedback integration
and personalized learning pathways.

\subsection{Backend Infrastructure}\label{backend-infrastructure}

The backend functions as the coordinating infrastructure of the system.
It manages user authentication, stores learning modules and dialogue
logs, and mediates real-time communication between frontend and AI
components. A scalable reactive architecture supports bidirectional
communication and modular expansion. Administrative management currently
operates via command-line tools; a graphical interface and REST API are
in development to facilitate browser-based configuration. Integration of
over 120 established MINT simulations is supported through an extensible
data model. Planned extensions include support for cooperative learning
scenarios and structured group work.

\subsection{Multi-Agent AI Subsystem}\label{multi-agent-ai-subsystem}

The AI subsystem represents the core innovation of AIDu. It consists of
specialized agents built on LLM architectures. Implemented agents
include an evaluation agent supporting analytical reasoning, a tutoring
agent for STEM domains, a language agent for conversational practice, an
exploration agent for interdisciplinary inquiry, a student simulation
agent for testing, and an experiential space agent coordinating
simulation-based learning.

Agents interact through structured tool calls, enabling them to update
simulation states, annotate the digital whiteboard, extract mathematical
expressions, perform concept analysis, or invoke helper agents.
Mathematical inputs can be formally validated using a Python-based
computer algebra system. For non-mathematical domains, automatic concept
graph generation models relationships among domain concepts and supports
structured feedback.

Currently, interaction flows are partially hard-coded. A dynamic AI
orchestrator is under development to select and combine agents
adaptively based on context and learner behavior. This orchestration
layer will enhance flexibility while preserving pedagogical control.

\subsection{Monitoring and Quality
Assurance}\label{monitoring-and-quality-assurance}

A defining characteristic of AIDu is real-time dialogue monitoring. Each
AI response is evaluated with respect to intended teaching behavior,
session objectives, and domain constraints. Speech-act-based analysis
allows detection of role violations, premature solution delivery, or
misinformation. This reliability layer is essential, as repeated
dialogue simulations demonstrate that identical prompts may otherwise
lead to divergent pedagogical trajectories.

Evaluation procedures include simulated learner testing, mathematical
solution validation, and structured discourse analysis. Planned
extensions involve live dialogue analytics providing teachers with
summaries of interaction depth, tutor behavior patterns, and learner
engagement metrics.

\section{Empirical Evaluation}\label{empirical-evaluation}

The system has been evaluated in multiple studies.

An exploratory study with eight students investigated usability and
dialogue patterns. Results revealed substantial variability in
interaction quality and informed refinements in monitoring mechanisms.

A larger study involving 117 students included structured usability
polling \citep{BuchnerSpahnEyer2026AIDuSGFB}
\citep{Langenegger2025LernenDurchKI} and qualitative analysis of 38
student--AI dialogues. Embedding dialogue within authentic environments
significantly reduced tutor role drift and stabilized epistemic support.

A controlled pretest--posttest study with 84 secondary school students
examined learning outcomes in a networking module focused on diagnosing
a malfunctioning printer. Results showed significant learning gains (d =
0.90). Notably, initial gender performance differences observed in the
pretest disappeared in the posttest, with particularly strong gains
among female students (d = 1.10). These findings suggest that epistemic
AI tutoring embedded in structured environments can both enhance
learning and contribute to reducing disparities in abstract STEM
domains.

Across studies, interaction patterns varied considerably. Some students
used the tutor reflectively for feedback, while others attempted to
extract full solutions. Without architectural safeguards, such pressures
risk undermining epistemic integrity. The monitoring layer proved
essential for maintaining pedagogical alignment.

\section{Contribution and Outlook}\label{contribution-and-outlook}

The AIDu project contributes a formalized model of epistemic AI
tutoring, a reliability-oriented system architecture, and empirical
evidence supporting competence-oriented AI integration in secondary
education. It demonstrates that embedding AI dialogue in authentic
contexts combined with real-time monitoring can stabilize tutor roles
and support reflective learning processes.

The findings underscore that educational impact depends less on raw
model capability than on architectural design. Reliable AI tutors
require explicit pedagogical framing, contextual grounding, dynamic
orchestration, and systematic quality assurance.

Future work will refine adaptive orchestration mechanisms, extend
longitudinal evaluation, and explore scalable deployment in diverse
educational settings. Further research will also investigate cooperative
AI-supported learning and automated real-time dialogue analytics.

\section{Conclusion}\label{conclusion}

AI integration in education is no longer optional. The decisive question
is whether systems are designed to promote durable understanding and
higher-order competencies. AIDu demonstrates that epistemic AI
tutoring---embedded in authentic learning environments and stabilized
through architectural monitoring---can support reflective, self-directed
learning while maintaining pedagogical integrity.

Design, not model size, determines educational reliability.


\input{EU_AI_ACT.tex}

% \clearpage
% \typeout{=== INPUTTING paper.bbl NOW ===}
% \input{paper.bbl}

\bibliographystyle{plainnat}
\bibliography{literature}

\end{document}

\bibliographystyle{plainnat}
\bibliography{literature}

\end{document}

\bibliographystyle{plainnat}
\bibliography{literature}

\end{document}

\bibliographystyle{plainnat}
\bibliography{literature}

\end{document}